\begingroup
\setlength{\intextsep}{2pt}
\captionsetup[lstlisting]{justification=raggedright,singlelinecheck=false}
\reservewrapspace{13\baselineskip}
\section{Motivating Example}
\label{sec:example}

\begin{wrapfigure}{r}{0.32\textwidth}
\begin{lstlisting}[style=pseudocode,caption={Two-pass source program.},label={lst:variance-before}]
$t\gets0$
for $j=0,\ldots,N-1$:
    $t\gets t+x_j$
$\mu\gets t/N$
$ss\gets0$
for $j=0,\ldots,N-1$:
    $ss\gets ss+(x_j-\mu)^2$
return $\;ss/N$
\end{lstlisting}
\end{wrapfigure}
方差展示了融合必须消除的依赖。
对于一行数据 $x_0,\ldots,x_{N-1}$，其中 $N>0$，\cref{lst:variance-before} 中的源程序先计算总和 $t$ 与均值 $\mu=t/N$，再累加中心化平方和 $ss$；方差为 $ss/N$。
只有第一次遍历得到最终均值后，第二次遍历才能开始原有计算。

为将这一依赖转化为合并操作，我们先考察使用局部均值计算的结果。
对于非空分块 $S$，令 $t_S$ 和 $ss_S$ 为源程序在该分块上的归约结果，$\mu_S$ 为该分块的均值。
直接计算这些结果仍需遍历该分块两次。
对于不相交分块 $\pa$ 和 $\pb$，沿用源归约运算（本例为 $\oplus=+$）得到 $ss_\pa\oplus ss_\pb=ss_\pa+ss_\pb$。
它通常不等于并集的中心化平方和，因为两个局部结果使用不同的均值，如 \cref{fig:variance}(a) 所示。
正确的合并操作必须处理均值变为共同均值时带来的变化。

因此，我们首先寻找无需共享均值即可组合的状态。
展开中心化归约，可以显式分离其对 $\mu_S$ 的依赖：
\begin{equation}
\label{eq:variance-reference}
ss_S=\sum_{j\in S}(x_j-\mu_S)^2
=\underbracket{\colorbox{boxgray}{$\displaystyle\sum_{j\in S}x_j^2$}}_{q_S}
 -2\mu_S\underbracket{\colorbox{boxgray}{$\displaystyle\sum_{j\in S}x_j$}}_{t_S}
 +\mu_S^2\underbracket{\colorbox{boxgray}{$\displaystyle\sum_{j\in S}1$}}_{n_S}.
\end{equation}
% Queue the overview with the reference-state discussion.
% Shared English visual; both prose editions include this source.
\par\medskip
\noindent\begin{minipage}{\linewidth}
\captionsetup{type=figure}
\centering
\begin{tikzpicture}[
  summary/.style={draw=rulegray,rounded corners=2pt,fill=boxgray!45,
    text width=3.75cm,minimum height=4.3cm,inner sep=5pt,
    align=center,font=\small},
  heading/.style={font=\sffamily\small\bfseries,anchor=south}]
\node[summary] (source) at (0,0) {
  $\displaystyle t=\sum_j x_j,\quad \mu=t/N$\\[5pt]
  $\displaystyle ss=\sum_j(x_j-\mu)^2$\\[9pt]
  Full row: $ss=\orig{\exTwoPass}$.\\[4pt]
  Adding local results:\\[2pt]
  $ss_\pa+ss_\pb=\stale{\exNaive}$
};
\node[summary] (reference) at (4.55,0) {
  Reference state tuple: $(t,q,n)$\\[5pt]
  $\begin{aligned}
    t&=\textstyle\sum_j x_j\\[2pt]
    q&=\textstyle\sum_j x_j^2\\[2pt]
    n&=\textstyle\sum_j 1
  \end{aligned}$\\[5pt]
  Componentwise join\\[3pt]
  $ss=q-t^2/n$\\[5pt]
  Shifted float32: $\stale{\exTextbookFloat}$
};
\node[summary] (maintained) at (9.1,0) {
  Maintained state tuple: $(t,ss,\aux{n})$\\[5pt]
  $\mu=t/n,\quad \mu_p=t_p/n_p$\\[5pt]
  $ss=ss_\pa+ss_\pb+\Delta$\\[3pt]
  $\displaystyle\Delta=\sum_{p\in\{\pa,\pb\}}n_p(\mu-\mu_p)^2$\\[7pt]
  Shifted float32: $\fix{\exTwoPassFloat}$
};
\node[heading,text=origblue] at ([yshift=3pt]source.north) {(a) Two-pass source};
\node[heading] at ([yshift=3pt]reference.north) {(b) Reference states};
\node[heading,text=fixgreen] at ([yshift=3pt]maintained.north) {(c) Maintained states};
\draw[flow] (source.east) -- (reference.west);
\draw[flow,draw=fixgreen] (reference.east) -- (maintained.west);
\end{tikzpicture}
\caption{Deriving a single-pass variance reducer.
For $\pa=(1,3)$ and $\pb=(2,6)$, adding locally centred results loses the change of mean.
The reference state tuple establishes a three-state bound; the maintained state tuple retains $t$ and $ss$ and adds \aux{$n$} for the correction.
The shifted results report $ss$ after adding $10^4$ to every input.}
\label{fig:variance}
% src: data/illus.json variance_exact and variance_float32
\end{minipage}\par\medskip

\reservewrapspace{10\baselineskip}
\begin{wrapfigure}{r}{0.32\textwidth}
\begin{lstlisting}[style=pseudocode,caption={Single pass with reference states.},label={lst:variance-reference}]
$(t,q,n)\gets(0,0,0)$
for $j=0,\ldots,N-1$:
    $t\gets t+x_j$
    $q\gets q+x_j^2$
    $n\gets n+1$
$ss\gets q-t^2/n$
return $\;ss/n$
\end{lstlisting}
\end{wrapfigure}
高亮的归约均不读取 $\mu_S$：它们分别是平方和 $q_S$、总和 $t_S$ 与计数 $n_S$。
它们构成子图 (b) 中的参考状态元组 $(t_S,q_S,n_S)$。
每个状态都可以通过加法组合，三者共同通过 $\mu_S=t_S/n_S$ 和 $ss_S=q_S-t_S^2/n_S$ 恢复源结果。
这就确立了一个状态数与输入长度无关的充分三状态表示；\cref{lst:variance-reference} 在同一次遍历中累积这些状态，再重构所需结果。

参考表示说明单遍计算在代数上可行，但它重构 $ss$ 时需要对两个可能很大且彼此接近的量做减法。
\cref{fig:variance} 中平移后的 float32 输入展示了由此产生的相消。
为保留中心化计算，\sys 转而从源结果 $(t,ss)$ 出发，分析修正这些结果需要哪些信息。
令 $\mu$ 为并集的均值，$\delta_\pa=\mu-\mu_\pa$ 为 $\pa$ 上的均值变化。
围绕局部均值展开可得
\begin{equation}
\label{eq:variance-rebase}
\sum_{j\in\pa}(x_j-\mu)^2
=ss_\pa-2\delta_\pa\sum_{j\in\pa}(x_j-\mu_\pa)+\delta_\pa^2 n_\pa
=ss_\pa+\delta_\pa^2 n_\pa.
\end{equation}
在实数语义下，由于相对于分块自身均值的偏差之和为零，一次项消失。
剩余修正项以 $n_\pa$ 加权，它表示中心发生变化的元素数量。
源结果 $(t_\pa,ss_\pa)$ 并未保留这一计数。
因此，上下文修正分析将 $n$ 识别为所选表示缺失的信息，\sys 随后将其作为辅助状态进行检验。
在本例中，仅这个系数就足以补全表示：由于 $n_S>0$，维护状态元组 $(t_S,ss_S,n_S)$ 可以恢复 $q_S=ss_S+t_S^2/n_S$，其余参考状态则可直接读取。
状态选择过程验证这一恢复关系后，才接受该元组。

\reservewrapspace{10\baselineskip}
\begin{wrapfigure}{r}{0.32\textwidth}
\begin{lstlisting}[style=pseudocode,caption={Single pass with maintained states.},label={lst:variance-after}]
$\leaf(x)=(x,0,1)$
$a\gets\leaf(x_0)$
for $j=1,\ldots,N-1$:
    $a\gets\Mrg(a,\leaf(x_j))$
$(t,ss,n)\gets a$
return $\;ss/n$
\end{lstlisting}
\end{wrapfigure}
上述上下文修正直接给出 \cref{fig:variance}(c) 中的操作 $\Mrg$：相加总和与计数，并在相加之前将每个局部中心化平方和修正到组合后的均值。
恢复关系提供了规定这一合并操作的另一种方式：恢复各分块的参考元组，相加参考分量，再重构并集的维护元组。
\sys 以这一构造为依据验证中心化公式，将选定状态与正确的合并操作联系起来，而执行时直接计算修正。
对于图示的平移输入的两个分块，这一修正合并与两遍程序的结果一致。

为将这一合并转化为融合，令每个分块恰好包含一个元素。
此时均值为 $x_j$，中心化平方和为零，计数为一，因此初始化函数 $\leaf(x_j)=(x_j,0,1)$ 可直接生成其单元素状态元组。
\Cref{lst:variance-after} 用 $a$ 保存累积元组，从第一个元素初始化，之后每读入一个输入元素，就将这一合并操作应用于相应的单元素状态元组。
这同时消除了分块内部的多次遍历和对全局均值的依赖：输入只读取一次，最终方差为 $ss/n$。

在实数语义下，所得合并恢复了 \citet{chan} 的成对方差公式；逐元素应用时，它在代数上等价于 Welford 更新~\citep{welford}。
同一并集性质还允许通过并行归约树组合。
\endgroup
